\chapter{The formal exponential and logarithm}\label{chap:weil-formal-exp-log}

This chapter is the power-series input for the zeta-function chapter: the
formal exponential $\operatorname{Exp}$, its inverse the formal logarithm
$\operatorname{Log}$, and the uniqueness principle for logarithmic
differential equations that produces both.  Nothing analytic happens: all
statements are coefficientwise identities of formal power series over a
commutative ring.

Throughout, $R$ is a commutative ring, $R[[T]]$ is the ring of formal power
series in one indeterminate (written \texttt{\detokenize{R⟦X⟧}} in Lean),
$c\colon R[[T]]\to R$ is the constant-coefficient ring homomorphism, and $D$
is the formal derivative $D(\sum a_nT^n)=\sum(n+1)a_{n+1}T^n$.  Mathlib
supplies $D$ as the derivation \texttt{\detokenize{PowerSeries.derivative}},
notation \texttt{\detokenize{d⁄dX}}, the exponential series
\texttt{\detokenize{PowerSeries.exp}}, the logarithmic series
\texttt{\detokenize{PowerSeries.log}}, and the substituted logarithm
\texttt{\detokenize{PowerSeries.logOf}}.  This chapter adds the missing
companion $\operatorname{Exp}$ and the inversion package relating the two.
The material is classical; it is the formal-series calculus underlying the
zeta function of a variety over a finite field, as in
\cite[Appendix~C]{Hartshorne}.

\paragraph{Coefficient-ring scope.}
Sections~\ref{sec:weil-log-ode} and~\ref{sec:weil-one-series} need only a
commutative ring, the first additionally torsion-free as an abelian group.
From Section~\ref{sec:weil-formal-exp} onwards $R$ is a commutative
$\mathbf Q$-algebra, and the inversion statements of
Section~\ref{sec:weil-exp-log-inverse} add torsion-freeness.  Each
declaration carries exactly the hypotheses it uses.

\section{Uniqueness for logarithmic differential equations}
\label{sec:weil-log-ode}

The whole chapter rests on one principle: a normalized power series is
determined by its logarithmic derivative.  It costs no denominators, so it is
stated over an arbitrary torsion-free commutative ring.  Recall that a power
series $A$ with $c(A)=1$ is a unit, with inverse
\texttt{\detokenize{PowerSeries.invOfUnit A 1}}, and that Mathlib's
\texttt{\detokenize{PowerSeries.derivative.ext}} says that $D$ together with
$c$ separates power series over a torsion-free ring.

\begin{lemma}[Normalized series with vanishing derivative]
  \label{lemma:weil-eq-one-of-derivative-eq-zero}
  \lean{PowerSeries.eq_one_of_derivative_eq_zero}
  \leanok

  Let $R$ be a torsion-free commutative ring and let $P\in R[[T]]$ satisfy
  $c(P)=1$ and $DP=0$.  Then $P=1$.
\end{lemma}

\begin{proof}
  \leanok
  $D1=0$ because $D$ is a derivation, and $c(1)=1$.  So $P$ and $1$ have the
  same derivative and the same constant coefficient, and
  \texttt{\detokenize{PowerSeries.derivative.ext}} applies.
\end{proof}

\begin{lemma}[Common logarithmic derivative]
  \label{lemma:weil-eq-of-derivative-eq-mul}
  \uses{lemma:weil-eq-one-of-derivative-eq-zero}
  \lean{PowerSeries.eq_of_derivative_eq_mul}
  \leanok

  Let $R$ be a torsion-free commutative ring and let $H,P,Q\in R[[T]]$
  satisfy
  \[
    c(P)=c(Q)=1,\qquad DP=HP,\qquad DQ=HQ .
  \]
  Then $P=Q$.
\end{lemma}

\begin{proof}
  \leanok
  Write $Q^{-1}$ for \texttt{\detokenize{invOfUnit Q 1}}, so that $QQ^{-1}=1$
  by \texttt{\detokenize{PowerSeries.mul_invOfUnit}}.  Differentiating that
  identity and using $DQ=HQ$ gives
  \[
    Q\,D(Q^{-1})+Q^{-1}HQ=0 ,
  \]
  and multiplying by $Q^{-1}$ gives $D(Q^{-1})=-HQ^{-1}$.  Hence, by the
  Leibniz rule,
  \[
    D\bigl(PQ^{-1}\bigr)
      = (DP)Q^{-1}+P\,D(Q^{-1})
      = HPQ^{-1}-PHQ^{-1}=0 .
  \]
  Moreover $c(PQ^{-1})=c(P)c(Q^{-1})=1$, because the constant coefficient of
  \texttt{\detokenize{invOfUnit Q 1}} is $1$.  So $PQ^{-1}=1$ by
  \Cref{lemma:weil-eq-one-of-derivative-eq-zero}, and multiplying by $Q$
  gives $P=Q$.
\end{proof}

\section{The submonoid of normalized series}
\label{sec:weil-one-series}

\begin{definition}[Normalized one-series]
  \label{def:weil-one-submonoid}
  \lean{PowerSeries.oneSubmonoid}
  \leanok

  Let $R$ be a commutative ring.  Define
  \[
    \operatorname{OneSeries}(R):=\{A\in R[[T]]\mid c(A)=1\},
  \]
  the kernel of $c$ viewed as a homomorphism of multiplicative monoids.  It
  is therefore a submonoid of $R[[T]]$ and carries the multiplication and the
  unit of $R[[T]]$.

  \formalizationnote{Dually, $\{U\mid c(U)=0\}$ is the ring-theoretic kernel
  $\ker c$, an ideal of $R[[T]]$, and is used in that form below; no separate
  definition is introduced for it.}
\end{definition}

\begin{lemma}[Membership in the one-series submonoid]
  \label{lemma:weil-mem-one-submonoid}
  \uses{def:weil-one-submonoid}
  \lean{PowerSeries.mem_oneSubmonoid}
  \leanok

  For $A\in R[[T]]$ one has $A\in\operatorname{OneSeries}(R)$ if and only if
  $c(A)=1$.
\end{lemma}

\begin{proof}
  \leanok
  This is the definition of the kernel of a monoid homomorphism.
\end{proof}

\section{The alternating geometric series}
\label{sec:weil-geometric}

\begin{lemma}[Inverse of $1+T$]
  \label{lemma:weil-mk-neg-one-pow-mul-one-add-X}
  \lean{PowerSeries.mk_neg_one_pow_mul_one_add_X}
  \leanok

  Let $R$ be a commutative ring.  Then
  \[
    \Bigl(\sum_{n\ge0}(-1)^nT^n\Bigr)\,(1+T)=1 .
  \]
\end{lemma}

\begin{proof}
  \leanok
  Compare coefficients.  In degree $0$ the left-hand side has coefficient
  $(-1)^0=1$.  In degree $n+1$ it has coefficient
  $(-1)^{n+1}+(-1)^n=0$.
\end{proof}

\begin{lemma}[Substituted derivative of the logarithm]
  \label{lemma:weil-subst-derivative-log}
  \uses{lemma:weil-mk-neg-one-pow-mul-one-add-X}
  \lean{PowerSeries.subst_derivative_log}
  \leanok

  Let $R$ be a commutative $\mathbf Q$-algebra and let $A\in R[[T]]$ satisfy
  $c(A)=1$.  Then
  \[
    \bigl(D\Lambda_R\bigr)\big|_{T=A-1}=A^{-1},
  \]
  where $\Lambda_R=\texttt{\detokenize{PowerSeries.log}}$ is the logarithmic
  series and $A^{-1}=\texttt{\detokenize{invOfUnit A 1}}$.
\end{lemma}

\begin{proof}
  \leanok
  Mathlib's \texttt{\detokenize{PowerSeries.deriv_log}} identifies
  $D\Lambda_R$ with $\sum_{n\ge0}(-1)^nT^n$.  Substitution at $A-1$, which is
  legitimate because $c(A-1)=0$, is a ring homomorphism, so applying it to
  \Cref{lemma:weil-mk-neg-one-pow-mul-one-add-X} and using
  $1+(A-1)=A$ gives $\bigl(D\Lambda_R\bigr)\big|_{T=A-1}\cdot A=1$.  A
  two-sided inverse in a commutative ring is unique, and $AA^{-1}=1$, so the
  two inverses of $A$ agree.
\end{proof}

\section{The formal exponential}
\label{sec:weil-formal-exp}

\begin{definition}[Formal exponential]
  \label{def:weil-formal-exp}
  \lean{PowerSeries.expOf}
  \leanok

  Let $R$ be a commutative $\mathbf Q$-algebra.  For $U\in R[[T]]$ define
  \[
    \operatorname{Exp}(U):=E_R\big|_{T=U},
    \qquad E_R:=\sum_{n\ge0}\frac{T^n}{n!},
  \]
  the substitution of $U$ into Mathlib's exponential series
  \texttt{\detokenize{PowerSeries.exp}}.

  \formalizationnote{\texttt{\detokenize{PowerSeries.subst}} is a total
  function, so $\operatorname{Exp}$ is defined on all of $R[[T]]$; the
  hypothesis $c(U)=0$, which makes the substitution meaningful, is carried by
  the individual lemmas.  The defining equation is
  \texttt{\detokenize{PowerSeries.expOf_eq}}.}
\end{definition}

\begin{lemma}[Constant coefficient of the exponential]
  \label{lemma:weil-constant-coeff-exp}
  \uses{def:weil-formal-exp}
  \lean{PowerSeries.constantCoeff_expOf}
  \leanok

  If $c(U)=0$, then $c(\operatorname{Exp}(U))=1$.
\end{lemma}

\begin{proof}
  \leanok
  Substitution at $U$ is a ring homomorphism, so
  $\operatorname{Exp}(U)-1=(E_R-1)\big|_{T=U}$.  The series $E_R-1$ has
  vanishing constant coefficient, and substituting a series with vanishing
  constant coefficient into one with vanishing constant coefficient again has
  vanishing constant coefficient
  (\texttt{\detokenize{PowerSeries.constantCoeff_subst_eq_zero}}).
\end{proof}

\begin{lemma}[Chain rule for the exponential]
  \label{lemma:weil-derivative-exp}
  \uses{def:weil-formal-exp}
  \lean{PowerSeries.derivative_expOf}
  \leanok

  If $c(U)=0$, then
  $D\operatorname{Exp}(U)=(DU)\operatorname{Exp}(U)$.
\end{lemma}

\begin{proof}
  \leanok
  By \texttt{\detokenize{PowerSeries.derivative_subst}}, $D(f|_{T=U})
  =(Df)|_{T=U}\cdot DU$; apply this to $f=E_R$ and use
  $DE_R=E_R$ (\texttt{\detokenize{PowerSeries.derivative_exp}}).
\end{proof}

\begin{lemma}[Chain rule for the logarithm]
  \label{lemma:weil-derivative-log}
  \uses{lemma:weil-subst-derivative-log}
  \lean{PowerSeries.derivative_logOf}
  \leanok

  Let $A\in R[[T]]$ satisfy $c(A)=1$.  Then
  \[
    D\operatorname{Log}(A)=(DA)A^{-1},
  \]
  where $\operatorname{Log}=\texttt{\detokenize{PowerSeries.logOf}}$ is
  Mathlib's substituted logarithm $\Lambda_R\big|_{T=A-1}$.
\end{lemma}

\begin{proof}
  \leanok
  Again by \texttt{\detokenize{PowerSeries.derivative_subst}},
  $D\operatorname{Log}(A)=\bigl(D\Lambda_R\bigr)\big|_{T=A-1}\cdot D(A-1)$.
  Now $D(A-1)=DA$, and the substituted factor is $A^{-1}$ by
  \Cref{lemma:weil-subst-derivative-log}.
\end{proof}

\section{Exponential and logarithm are inverse}
\label{sec:weil-exp-log-inverse}

From here on $R$ is a torsion-free commutative $\mathbf Q$-algebra.

\begin{lemma}[Exponential of a sum]
  \label{lemma:weil-exp-add}
  \uses{lemma:weil-eq-of-derivative-eq-mul,lemma:weil-constant-coeff-exp,
    lemma:weil-derivative-exp}
  \lean{PowerSeries.expOf_add}
  \leanok

  If $c(U)=c(V)=0$, then
  $\operatorname{Exp}(U+V)=\operatorname{Exp}(U)\operatorname{Exp}(V)$.
\end{lemma}

\begin{proof}
  \leanok
  Both sides have constant coefficient $1$ and both solve the logarithmic
  differential equation with $H=DU+DV$: for the left-hand side this is
  \Cref{lemma:weil-derivative-exp} applied to $U+V$, and for the right-hand
  side it is the Leibniz rule together with
  \Cref{lemma:weil-derivative-exp} for $U$ and for $V$.  Conclude by
  \Cref{lemma:weil-eq-of-derivative-eq-mul}.
\end{proof}

\begin{lemma}[Exponential of zero]
  \label{lemma:weil-exp-zero}
  \uses{lemma:weil-eq-one-of-derivative-eq-zero,lemma:weil-derivative-exp}
  \lean{PowerSeries.expOf_zero}
  \leanok

  $\operatorname{Exp}(0)=1$.
\end{lemma}

\begin{proof}
  \leanok
  $c(\operatorname{Exp}(0))=1$ and
  $D\operatorname{Exp}(0)=(D0)\operatorname{Exp}(0)=0$, so
  \Cref{lemma:weil-eq-one-of-derivative-eq-zero} applies.
\end{proof}

\begin{proposition}[Logarithm after exponential]
  \label{prop:weil-log-exp}
  \uses{lemma:weil-derivative-log,lemma:weil-derivative-exp,
    lemma:weil-constant-coeff-exp}
  \lean{PowerSeries.logOf_expOf}
  \leanok

  If $c(U)=0$, then $\operatorname{Log}(\operatorname{Exp}(U))=U$.
\end{proposition}

\begin{proof}
  \leanok
  Write $A=\operatorname{Exp}(U)$, so $c(A)=1$.  By
  \Cref{lemma:weil-derivative-log} and \Cref{lemma:weil-derivative-exp},
  \[
    D\operatorname{Log}(A)=(DA)A^{-1}=(DU)AA^{-1}=DU .
  \]
  Also $c(\operatorname{Log}(A))=0=c(U)$ by
  \texttt{\detokenize{PowerSeries.constantCoeff_logOf}}.  Conclude with
  \texttt{\detokenize{PowerSeries.derivative.ext}}.
\end{proof}

\begin{proposition}[Exponential after logarithm]
  \label{prop:weil-exp-log}
  \uses{lemma:weil-eq-of-derivative-eq-mul,lemma:weil-derivative-log,
    lemma:weil-derivative-exp}
  \lean{PowerSeries.expOf_logOf}
  \leanok

  If $c(A)=1$, then $\operatorname{Exp}(\operatorname{Log}(A))=A$.
\end{proposition}

\begin{proof}
  \leanok
  Put $H=(DA)A^{-1}$.  Both $\operatorname{Exp}(\operatorname{Log}(A))$ and
  $A$ have constant coefficient $1$.  The first satisfies
  $D\operatorname{Exp}(\operatorname{Log}(A))
   =\bigl(D\operatorname{Log}(A)\bigr)\operatorname{Exp}(\operatorname{Log}(A))
   =H\operatorname{Exp}(\operatorname{Log}(A))$ by
  \Cref{lemma:weil-derivative-exp} and \Cref{lemma:weil-derivative-log}, and
  the second satisfies $DA=(DA)A^{-1}A=HA$.  Conclude by
  \Cref{lemma:weil-eq-of-derivative-eq-mul}.
\end{proof}

\begin{lemma}[Logarithm of one]
  \label{lemma:weil-log-one}
  \uses{prop:weil-log-exp,lemma:weil-exp-zero}
  \lean{PowerSeries.logOf_one}
  \leanok

  $\operatorname{Log}(1)=0$.
\end{lemma}

\begin{proof}
  \leanok
  $\operatorname{Log}(1)=\operatorname{Log}(\operatorname{Exp}(0))=0$ by
  \Cref{lemma:weil-exp-zero} and \Cref{prop:weil-log-exp}.
\end{proof}

\begin{proposition}[Logarithm of a product]
  \label{prop:weil-log-mul}
  \uses{lemma:weil-exp-add,prop:weil-log-exp,prop:weil-exp-log}
  \lean{PowerSeries.logOf_mul}
  \leanok

  If $c(A)=c(B)=1$, then
  $\operatorname{Log}(AB)=\operatorname{Log}(A)+\operatorname{Log}(B)$.
\end{proposition}

\begin{proof}
  \leanok
  By \Cref{prop:weil-exp-log} and \Cref{lemma:weil-exp-add},
  \[
    AB=\operatorname{Exp}(\operatorname{Log}(A))
       \operatorname{Exp}(\operatorname{Log}(B))
      =\operatorname{Exp}\bigl(\operatorname{Log}(A)+\operatorname{Log}(B)\bigr).
  \]
  Both logarithms have vanishing constant coefficient, so applying
  \Cref{prop:weil-log-exp} to the right-hand side gives the claim.
\end{proof}

\begin{proposition}[Logarithm of a finite product]
  \label{prop:weil-log-prod}
  \uses{prop:weil-log-mul,lemma:weil-log-one}
  \lean{PowerSeries.logOf_prod}
  \leanok

  Let $S$ be a finite index set and let $A_i\in R[[T]]$ satisfy $c(A_i)=1$
  for every $i\in S$.  Then
  \[
    \operatorname{Log}\Bigl(\prod_{i\in S}A_i\Bigr)
      =\sum_{i\in S}\operatorname{Log}(A_i).
  \]
\end{proposition}

\begin{proof}
  \leanok
  Induction on $S$.  The empty case is \Cref{lemma:weil-log-one}.  For the
  induction step, the product over the smaller set again has constant
  coefficient $1$, so \Cref{prop:weil-log-mul} applies.
\end{proof}

\begin{theorem}[Exponential/logarithm bijection]
  \label{thm:weil-exp-log-equiv}
  \uses{def:weil-one-submonoid,lemma:weil-exp-add,prop:weil-log-exp,
    prop:weil-exp-log,lemma:weil-constant-coeff-exp}
  \lean{PowerSeries.expOfMulEquiv}
  \leanok

  Let $R$ be a torsion-free commutative $\mathbf Q$-algebra.  The formal
  exponential is an isomorphism
  \[
    \operatorname{Exp}\colon
      \bigl(\ker c,+\bigr)\;\xrightarrow{\ \sim\ }\;
      \bigl(\operatorname{OneSeries}(R),\cdot\bigr)
  \]
  from the additive group of the ideal $\ker c\subseteq R[[T]]$ onto the
  multiplicative monoid of normalized one-series, with inverse
  $\operatorname{Log}$.
\end{theorem}

\begin{proof}
  \leanok
  \Cref{lemma:weil-constant-coeff-exp} and
  \texttt{\detokenize{PowerSeries.constantCoeff_logOf}} show that
  $\operatorname{Exp}$ and $\operatorname{Log}$ restrict to maps between the
  two carriers; \Cref{prop:weil-log-exp} and \Cref{prop:weil-exp-log} show
  that they are mutually inverse; \Cref{lemma:weil-exp-add} is
  multiplicativity.

  \formalizationnote{In Lean the source is written
  \texttt{\detokenize{Multiplicative (RingHom.ker (constantCoeff (R := R)))}}
  so that the statement is a \texttt{\detokenize{MulEquiv}}; the projection
  lemmas are \texttt{\detokenize{PowerSeries.expOfMulEquiv_apply}} and
  \texttt{\detokenize{PowerSeries.expOfMulEquiv_symm_apply}}.}
\end{proof}
